\section{Inverse kinematics and the capture manoeuvre}
\label{sec:ik}

\subsection{A five-joint arm in a six-dimensional task space}

The manipulator has five joints; a rigid-body pose has six degrees of freedom.
The pose problem is therefore structurally unsolvable in general, and a solver
that hides this fact by converging to something is worse than one that does not
converge. The approach taken here states the deficiency in the weighting and
reports the residual it leaves.

\subsection{Spatial error}

Given a target pose $(\mathbf p_t, R_t)$ and the current pose
$(\mathbf p_k, R_k)$ with columns $\mathbf x_k, \mathbf y_k, \mathbf z_k$, the
task error is

\begin{equation}
\Delta\mathbf p = \mathbf p_t - \mathbf p_k, \qquad
\mathbf e_o = \tfrac12\left(\mathbf x_k\times\mathbf x_t + \mathbf y_k\times\mathbf y_t + \mathbf z_k\times\mathbf z_t\right),
\qquad
\Delta\mathbf x = \begin{bmatrix}\Delta\mathbf p \\ \mathbf e_o\end{bmatrix}.
\label{eq:taskerr}
\end{equation}

The orientation error uses the Siciliano form \cite{siciliano} rather than
axis-angle extraction, which is singular at $\pm180^\circ$. That is not a
theoretical concern here: the seeds used are far from the target and rotations
near $180^\circ$ genuinely occur on the first iterations.

\subsection{Weighted damped least squares}

The geometric Jacobian is assembled column by column for revolute joints,

\begin{equation}
J_i = \begin{bmatrix} \mathbf z_{i}\times(\mathbf p_{EE} - \mathbf r_{i}) \\ \mathbf z_{i}\end{bmatrix},
\label{eq:jacobian}
\end{equation}

ordered $[\text{linear};\text{angular}]$ to match~\eqref{eq:taskerr}, and the
update is

\begin{equation}
\boldsymbol\theta \leftarrow \boldsymbol\theta + \alpha\left(J^\top W J + \lambda^2 I\right)^{-1} J^\top W \Delta\mathbf x,
\label{eq:wdls}
\end{equation}

with $\lambda = 0.05$, $\alpha = 0.4$ and
$W = \mathrm{diag}(1,1,1,1,1,0.01)$. The damping keeps the update bounded near
singular configurations; the weighting spends the available freedom on position
and de-emphasises the one attitude component a five-joint wrist cannot serve.

\subsection{A real singularity at the obvious seed}

The stowed configuration $\boldsymbol\theta = \mathbf 0$ is not a neutral
starting point. Fully extended along $z$, joint~1 rotates about an axis passing
through the end effector, so it produces no end-effector translation at all and
the position Jacobian loses a column. The solver escapes because the attitude
term still torques joint~1, but the code carries a deterministic seed cascade
anyway and logs which seed was used, rather than pretending the singularity is
not there.

\subsection{Results}

Two targets are run. The first is the forward-kinematics pose of a known
configuration, so an exact solution provably exists and the convergence claim
means something. Starting from the zero pose the solver reaches
$\mIKres\unit{mm}$ in $\mIKiters$ iterations, well inside the $80$-iteration
budget.

The solution found is \emph{not} the configuration the target was generated
from: the solver returns $[40.0, 16.3, 20.1, 28.6, 25.0]^\circ$ against the
generating $[40, 35, -20, 50, 25]^\circ$. Both reach the same pose because
joints 2, 3 and 4 rotate about a common axis and only their sum enters -- $65^\circ$
in both cases. Multiple inverse solutions are the expected behaviour of a serial
chain, not a defect.

The second target is operationally motivated: a point $0.4\unit{m}$ in front of
the mothership docking face with the tool axis along $+x$. Position converges to
$\mIKtwoRes\unit{mm}$ and a small attitude residual survives, which is precisely
the structural deficiency admitted at the start of this section.

\begin{figure}[htbp]
\centering
\includegraphics[width=\textwidth]{fig20_ik_convergence.png}
\caption{Left: WDLS convergence for both targets, log scale. Position residual
falls below $1\unit{mm}$ in $\mIKiters$ iterations. Right: the capture path of
the end effector from the stowed pose to the solution, with the final
configuration overlaid.}
\label{fig:ik}
\end{figure}

\subsection{Rest-to-rest joint trajectory}

The capture is flown as a cubic in joint space over $t_f = 30\unit{s}$ with zero
rate at both ends,

\begin{equation}
\theta(t) = \theta_0 + \frac{3\Delta\theta}{t_f^2}t^2 - \frac{2\Delta\theta}{t_f^3}t^3,
\qquad
\dot\theta(t) = \frac{6\Delta\theta}{t_f^2}t - \frac{6\Delta\theta}{t_f^3}t^2 .
\label{eq:cubic}
\end{equation}

A cubic rather than a quintic is a deliberate choice: berthing rates are slow
enough that the discontinuous acceleration at the endpoints is harmless, and the
velocity profile stays trivially verifiable. The peak rate is
$1.5\,\Delta\theta/t_f$ at mid-time, here $0.035\unit{rad/s}$, and the end
effector travels $4.12\unit{m}$ along the capture path.

\begin{figure}[htbp]
\centering
\includegraphics[width=0.95\textwidth]{fig17_joint_trajectory.png}
\caption{Rest-to-rest cubic joint trajectory over $30\unit{s}$. Rates start and
end at exactly zero, which is checked to $10^{-17}\unit{rad/s}$ rather than read
off the plot.}
\label{fig:cubic}
\end{figure}

\subsection{Verification}

The boundary conditions of~\eqref{eq:cubic} are checked symbolically against the
sampled trajectory: $\theta(0) = \theta_0$ and $\theta(t_f) = \theta^*$ to
$10^{-16}$, $\dot\theta(0) = \dot\theta(t_f) = 0$ to $10^{-17}$, and the peak
rate matches the analytical $1.5\Delta\theta/t_f$ exactly. The inverse-kinematics
solution is verified the only way that is meaningful for an iterative solver: by
pushing the returned joint vector back through the forward kinematics of
Section~\ref{sec:robotics} and measuring the pose it actually produces.
